% ------------------------------------------------------------------------------
% 3. OBJETIVOS
% ------------------------------------------------------------------------------
\subsection{3. Objetivos}\label{sec:objetivos}

\subsubsection{3.1. Objetivo General}\label{sec:objetivo_general}

Desarrollar un Sistema de Información Web y Móvil para la Postulación y Preselección Automatizada de la Beca Comedor mediante Inteligencia Artificial, con Reportes de Business Intelligence para la DUBSS - UAGRM, que permita automatizar la postulación en línea, precalificar algorítmicamente la vulnerabilidad socioeconómica de los estudiantes, gestionar la asistencia física/digital y proveer dashboards analíticos y proyecciones predictivas que apoyen la toma de decisiones institucionales y auditorías.

\subsubsection{3.2. Objetivos Específicos}\label{sec:objetivos_especificos}

\begin{itemize}
    \item Capturar y determinar los requisitos funcionales y no funcionales del proceso de gestión del comedor universitario mediante entrevistas, revisión del flujo de postulación masiva (más de 20 000 carpetas), normativas de auditoría física (firmas y carnet) y los flujos operativos de la DUBS - UAGRM.

    \item Analizar y modelar los procesos del sistema, definiendo la arquitectura del software web y móvil, las reglas de negocio para la puntuación de vulnerabilidad, las variables clave para los módulos de Business Intelligence y los requerimientos de datos para los modelos de Machine Learning predictivo.

    \item Diseñar la arquitectura lógica, física y de interfaces del sistema, incluyendo el diseño de la base de datos, maquetación de las interfaces para usuarios (plataforma web y app móvil), el modelo algorítmico de priorización de carpetas, las vistas de reportes BI y la estructura del modelo predictivo.

    \item Implementar el sistema de información web y móvil, codificando los módulos de postulación en línea, la precalificación automática de solicitudes, la gestión administrativa de los 900 becados, el registro e impresión de planillas de asistencia para auditoría, los reportes interactivos de Business Intelligence y los algoritmos de Machine Learning Predictivo.

    \item Realizar las pruebas del sistema (unitarias, de integración, de carga para la atención masiva de carpetas y de usabilidad en entornos web y móvil), validando el correcto funcionamiento de los algoritmos predictivos, los reportes BI y garantizando la seguridad e integridad de la información socioeconómica.
\end{itemize}
